% CO-AUTHOR EDIT: INSERT RATIONALITY TEST [T2.1]
\begin{table}[htbp]
\centering
\caption{FL Participation Rationality Test}
\label{tab:rationality}
\begin{threeparttable}
\small
\begin{tabular}{lcccccc}
\toprule
 & $N$ & Mean & Win & $E[\pi]$ & $\%$ negative & Career \\
 & firms & particip. & rate & per bid & $E[\pi]$ & loss (R\$K) \\
\midrule
FL firms & 2,444 & 27.7 & 0.000 & -755 & 100.0\% & 10.6 \\
Non-FL always-losers & 11,762 & 4.1 & 0.000 & -2,390 & 100.0\% & 5.3 \\
Low win-rate ($<$10\%) & 5,476 & 618.2 & 0.055 & 399 & 23.5\% & --- \\
Regular competitors & 16,367 & 537.8 & 0.396 & 6,801 & 0.3\% & --- \\
\bottomrule
\end{tabular}
\begin{tablenotes}
\small
\item \textit{Notes:} Expected profit per bid assumes 10\% gross margin on wins
and 1\% of median tender value as bidding cost.
$E[\pi] = \text{win\_rate} \times \text{mean\_value} \times 0.10 - 0.01 \times \text{median\_value}$.
Career loss = bidding cost $\times$ total participations.
FL firms have zero wins by construction; under competitive bidding,
their participation generates guaranteed losses for any positive bidding cost.
\end{tablenotes}
\end{threeparttable}
\end{table}
